% =====================================================================
%  comandos.tex  ->  caixas didaticas e placeholder de figura
% =====================================================================

% ---------- Caixa: Ideia-chave ---------------------------------------
\newtcolorbox{ideiachave}[1][Ideia-chave]{
  enhanced, breakable,
  colback=btFundo, colframe=btDecisao,
  fonttitle=\bfseries, title=#1,
  left=10pt, right=10pt, top=8pt, bottom=8pt,
  boxrule=0.8pt, arc=3pt
}

% ---------- Caixa: Passo a passo -------------------------------------
\newtcolorbox{passos}[1][Passo a passo]{
  enhanced, breakable,
  colback=white, colframe=btValido,
  fonttitle=\bfseries, title=#1,
  left=10pt, right=10pt, top=8pt, bottom=8pt,
  boxrule=0.8pt, arc=3pt
}

% ---------- Caixa: Atencao / limitacao -------------------------------
\newtcolorbox{atencao}[1][Atencao]{
  enhanced, breakable,
  colback=white, colframe=btPodado,
  fonttitle=\bfseries, title=#1,
  left=10pt, right=10pt, top=8pt, bottom=8pt,
  boxrule=0.8pt, arc=3pt
}

% ---------- Caixa: Quadro-resumo executivo ---------------------------
\newtcolorbox{quadroresumo}[1][Em resumo]{
  enhanced, breakable,
  colback=btFundo, colframe=btNeutro,
  fonttitle=\bfseries\color{btTexto}, title=#1,
  left=10pt, right=10pt, top=8pt, bottom=8pt,
  boxrule=0.8pt, arc=3pt
}

% ---------------------------------------------------------------------
%  \figura{rotulo}{legenda}{caminho do arquivo TikZ (sem extensao)}
%  Insere uma figura TikZ finalizada, centralizada, com legenda e label.
% ---------------------------------------------------------------------
\newcommand{\figura}[3]{%
  \begin{figure}[H]
    \centering
    \input{#3}
    \caption{#2}
    \label{#1}
  \end{figure}%
}

% ---------------------------------------------------------------------
%  \figuraplaceholder{rotulo}{legenda}{descricao da figura desejada}
%  Reserva o espaco de uma figura ainda nao produzida, mostrando uma
%  caixa tracejada com a descricao do que ela deve conter. Mantem a
%  numeracao e a legenda corretas (\ref e lista de figuras funcionam).
% ---------------------------------------------------------------------
\newcommand{\figuraplaceholder}[3]{%
  \begin{figure}[H]
    \centering
    \begin{tcolorbox}[
        enhanced, width=0.85\linewidth, halign=center, valign=center,
        colback=btFundo, frame hidden,
        borderline={0.8pt}{0pt}{btNeutro,dashed},
        arc=4pt, height=4.5cm
      ]
      {\bfseries\color{btNeutro} [ FIGURA A PRODUZIR ]}\\[6pt]
      {\small\itshape #3}
    \end{tcolorbox}
    \caption{#2}
    \label{#1}
  \end{figure}%
}
